\section{Một trò chơi chia đóng góp}
\label{sec:ch04-tro-choi-chia-gia-tri}

\index{trò chơi hợp tác}%
Trước khi xét \tn{phân phối nền}{background distribution} của SHAP, ta cần xác
định quy tắc mà nó dùng để
chia đóng góp. Quy tắc đó bắt đầu từ bài toán gốc của lý thuyết trò chơi hợp
tác. Gọi $N$ là tập các người chơi và gọi $v(S)$ là giá trị mà liên minh
$S \subseteq N$ tạo ra; ký hiệu $F$ được phụ lục~\ref{app:ky-hieu} dành cho
hàm điểm số đang được attribution, nên trò chơi dùng chữ khác.
Một người chơi không có một đóng góp cố định: lợi ích khi thêm nó vào liên
minh phụ thuộc những người chơi đã có mặt. \tn{Giá trị Shapley}{Shapley value}
chia tổng giá trị bằng cách lấy trung bình đóng góp biên của mỗi người chơi
trên mọi thứ tự tham gia có thể có.

Với một \tn{đặc trưng}{feature} $i$, phương trình là

\begin{equation}
\label{eq:ch04-gia-tri-shapley}
\phi_i(v)=
\sum_{S \subseteq N \setminus \{i\}}
\frac{|S|!(|N|-|S|-1)!}{|N|!}
\left[v(S\cup\{i\})-v(S)\right].
\end{equation}

Hệ số phía trước ngoặc vuông là xác suất để, trong một hoán vị ngẫu nhiên của
$N$, đúng các phần tử của $S$ đứng trước $i$. Vì vậy,
phương trình~\ref{eq:ch04-gia-tri-shapley} không ưu ái thứ tự nào. Nó chỉ yêu cầu ta biết
giá trị thay đổi bao nhiêu khi $i$ được thêm vào sau mỗi liên minh có thể có.

\index{giá trị Shapley|(}%
Kết quả cổ điển được đặc trưng bằng bốn tiên đề: hiệu quả, đối xứng, người
chơi rỗng và \tn{tính tuyến tính}{linearity}. Hiệu quả yêu cầu tổng phần chia bằng giá trị của
toàn liên minh trừ giá trị của liên minh rỗng. Đối xứng buộc hai người chơi
luôn có cùng đóng góp biên phải nhận cùng giá trị. Người chơi rỗng nhận giá
trị không vì không làm thay đổi liên minh nào. Tính tuyến tính giữ cách chia
nhất quán khi cộng hoặc nhân các trò chơi với hằng số~\autocite{p19shapinference}.
